\documentclass[11pt, a4paper]{article}

% --- Paquets ---
\usepackage[utf8]{inputenc}
\usepackage[T1]{fontenc}
\usepackage[french]{babel}
\usepackage[margin=2.5cm]{geometry}
\usepackage{hyperref}
\usepackage{enumitem}
\usepackage{booktabs}
\usepackage{listings}
\usepackage{xcolor}

% --- Configuration Code ---
\lstset{
    language=Python,
    basicstyle=\ttfamily\small,
    keywordstyle=\color{blue},
    stringstyle=\color{red},
    commentstyle=\color{gray},
    breaklines=true,
    frame=single,
}

\title{\textbf{Rapport de recherche : Bibliothèque ByzFL} \\ \large Étude de la robustesse dans l'apprentissage fédéré}
\author{ABBES Samy \& AREZKI Célia}
\date{Semaine de 04 Mai 2026}

\begin{document}

\maketitle

\section{Introduction}
Ce rapport présente le travail effectué lors de notre première phase de recherche sur l'apprentissage fédéré avec \textbf{ByzFL}. Nous y détaillons notre processus d'apprentissage de l'outil \textbf{ByzFL}, la conception de notre première expérience sur le jeu de données MNIST, ainsi que l'évolution de nos résultats suite aux corrections techniques apportées.

\section{Prise en main de la bibliothèque}
La première semaine de notre TER a été consacrée à l'exploration et à la prise en main de la bibliothèque \textbf{ByzFL}. Cette phase nous a permis de comprendre l'implémentation modulaire de l'outil à travers ses quatre composantes fondamentales :

\begin{itemize}
    \item \textbf{DataDistributor} : Responsable de la partition des données entre les clients (distributions IID ou non-IID).
    \item \textbf{Client} : Simule les entités honnêtes réalisant l'entraînement local et le calcul des gradients.
    \item \textbf{ByzantineClient} : Simule les attaquants capables de corrompre les gradients avant leur envoi.
    \item \textbf{Server} : Point central qui orchestre l'agrégation des gradients reçus et met à jour le modèle global.
\end{itemize}

Nous avons également étudié les différentes méthodes d'agrégation robuste (\textit{TrMean, Krum, Median...etc}) et les types d'attaques (\textit{SignFlipping, IPM, ALIE...etc}) proposés par la bibliothèque pour comprendre leurs fondements mathématiques.

\section{Contenu de l'expérience MNIST}
Notre première expérience pratique vise à évaluer la résistance d'un modèle de réseau de neurones (\textit{cnn\_mnist}) face à diverses perturbations. Le but est de confronter des mécanismes de défense à des stratégies d'attaque spécifiques.

L'expérience consiste à comparer plusieurs combinaisons :
\begin{itemize}
    \item \textbf{Baseline} : Une agrégation simple (\textit{Average}) sans aucune attaque pour établir le score de référence.
    \item \textbf{Attaque brute vs Défense simple} : \textit{SignFlipping} contre une moyenne (\textit{Average}) et contre une médiane (\textit{Median}).
    \item \textbf{Attaque optimisée vs Défense statistique} : \textit{Inner Product Manipulation (IPM)} contre la moyenne tronquée (\textit{TrMean}).
    \item \textbf{Attaque furtive vs Défense géométrique} : \textit{A Little Is Enough (ALIE)} contre l'algorithme \textit{Krum}.
\end{itemize}

L'objectif global est d'observer si les agrégateurs robustes parviennent à maintenir une précision élevée là où une moyenne classique échouerait.

\section{Historique et Évolution de l'expérience}

\subsection{Phase 1 : Résultats initiaux et difficultés}
Lors de nos premiers tests, nous avons obtenu des résultats très médiocres. Le modèle ne parvenait pas à apprendre, avec une précision finale stagnante aux alentours de 10-13\%, ce qui correspondait quasiment à un choix aléatoire sur les 10 classes de MNIST.

\textbf{Paramétrages de l'époque :}
\begin{itemize}
    \item \textbf{Nombre de rounds} : 3 (trop faible pour une convergence).
    \item \textbf{Momentum} : 0.9 (potentiellement trop instable pour ce début d'entraînement).
\end{itemize}

\subsection{Phase 2 : Correction et succès}
Après une vérification rigoureuse du code avec notre professeur, nous avons identifié que le nombre de cycles d'apprentissage était insuffisant et que certains hyperparamètres nuisaient à la stabilité. En corrigeant ces éléments, nous avons obtenu des résultats conformes aux attentes.

\textbf{Nouveaux paramétrages :}
\begin{itemize}
    \item \textbf{Nombre de rounds} : 1000 (permettant une convergence complète).
    \item \textbf{Momentum} : 0.1 (pour une mise à jour plus stable des gradients).
\end{itemize}

\section{Résultats finaux (Distribution IID)}
Voici les performances obtenues après correction du protocole :

\section{Résultats finaux et Comparaison}
Le tableau ci-dessous présente la comparaison entre notre première tentative (infructueuse) et l'expérience corrigée. On observe que l'ajustement du nombre de rounds et du momentum a permis au modèle de passer d'un état de non-apprentissage à une convergence quasi-parfaite.

\begin{table}[h]
\centering
\begin{tabular}{@{}lllcc@{}}
\toprule
\textbf{Scénario} & \textbf{Agrégateur} & \textbf{Attaque} & \textbf{Acc. (3 rounds)} & \textbf{Acc. (1000 rounds)} \\ \midrule
clean\_average     & Average            & None             & 0.1217                  & \textbf{0.9876}            \\
signflip\_average  & Average            & SignFlipping     & 0.1202                  & 0.9849                     \\
signflip\_median   & Median             & SignFlipping     & 0.1161                  & 0.9832                     \\
ipm\_trmean        & TrMean             & IPM              & 0.1305                  & 0.9777                     \\
alie\_krum         & Krum               & ALIE             & 0.1041                  & 0.9831                     \\ \bottomrule
\end{tabular}
\caption{Comparaison des performances : Expérience initiale vs Expérience corrigée}
\end{table}

\section{Analyse et Conclusion}
Le passage de 3 à 1000 rounds a transformé une expérience médiocre en un succès technique. Nous avons validé que dans une distribution IID avec un seul attaquant, les défenses proposées par \textbf{ByzFL} sont extrêmement efficaces, limitant la perte de précision à moins de 1\% par rapport au cas sain. La prochaine étape consistera à complexifier la distribution des données (Non-IID) pour tester les limites de ces agrégateurs.

\end{document}
